%% Ejemplo de la plantilla de latex para las diapositivas del Máster en IA 

%%%%%%%%%%%%%%%%%%%%%%%%%%%%
%% Valores para el ratio de aspecto: 43, 169 
%\documentclass[10pt, envcountsect, presentation, aspectratio=1610, hyperref={hypertexnames=true}]{beamer}
\documentclass[10pt, envcountsect, handout, aspectratio=1610, hyperref={hypertexnames=true}]{beamer}

%%%%%%%%%%%%%%%%%%%%%%%%%%%%

%%%%%%%%%%%%%%%%%%%%%%%%%%%%
%% Incluir fichero de estilo del máster
\input ../style/plantillaMasterIA.sty
%%%%%%%%%%%%%%%%%%%%%%%%%%%%

\addbibresource{../references.bib}



%\setbeameroption{show only notes} % Only notes
%\setbeameroption{show notes on second screen=right}  % Both
%\setbeamertemplate{note page}[compress]
\setbeameroption{hide notes} % Only slides



\renewcommand*\footnoterule{}
\setbeamerfont{footnote}{size=\tiny}




%\input{embed_video.tex}
%\usepackage{movie15}b
\usepackage{multimedia}
%

%%%%%%%%%%%%%%%%%%%%%%%%%%%%
%% Información que aparecerá en la portada
\title[Aprendizaje por Refuerzo]{Aprendizaje por Refuerzo}
\subtitle{Métodos de n-Pasos} % short title for footer

\author[L. D. Hernández] % Para el pie de página, poner los autores abreviados separados por comas.
{%
	\sc{Luis Daniel Hernández Molinero }\\ % Autor 1
	\textit{Ciencia de la Computación e Inteligencia Artificial}\\  % Area de conocimiento
	\textit{Dept. Ingeniería de la Información y las Comunicaciones}\\  % Departamento
	\\   % No dejes espacio entre esta línea y la llave siguiente 
}

\institute[DIIC]% % Poner las iniciales de los diferentes departamentos de los autores separadas por comas.
{
	\textit{Universidad de Murcia}
}


\newcounter{lastyear}
\setcounter{lastyear}{\the\year}
\addtocounter{lastyear}{-1}

\date{\arabic{lastyear} / \the\year} % Curso académico

%%%%%%%%%%%%%%%%%%%%%%%%%%%%%%%%%%
%% Contenido de las diapositivas
\setbeamerfont{normal text}{size=\normalsize} % Modifica el tamaño del texto de las diapositivas.
\AtBeginDocument{\usebeamerfont{normal text}}

%%%%%%%%%%%%%%%%%%%%%%%%%%%%%%%%%%

\begin{document}	

%%%%%%%%%%%%%%%%%%%%%%%%%%%%%%%%%%









%%%%%%%%%%%%%%%%%%%%%%%%%%%%%%%%%%

\begin{frame}[plain]
	\titlepage
\end{frame}

%%%%%%%%%%%%%%%%%%%%%%%%%%%%%%%%%%

\begin{frame}{Índice de Contenidos}
	\tableofcontents 
	
\note{

}

\end{frame}








%%%%%%%%%%%%%%%%%%%%%%%%%%%%%%%%%%
%%%%%%%%%%%%%%%%%%%%%%%%%%%%%%%%%%
\section{Introducción}

%%%%%%%%%%%%%%%%%%%%%%%%%%%%%%%%%%
%%%%%%%%%%%%%%%%%%%%%%%%%%%%%%%%%%


%%%%%%%%%%%%%%%%%%%%%%%%%%%%%%%%%%
%%%%%%%%%%%%%%%%%%%%%%%%%%%%%%%%%%
\begin{frame}{Introducción}
\begin{itemize}
	\item Un \key{resumen} de los métodos MC y TD(0) Hemos estudiado:
	
	\centerline{\scalebox{.8}{
		\begin{minipage}{1\textwidth}
		\begin{tabularx}{\textwidth}{|l|Y|Y|}
			\hline
			\textbf{Característica}        & \textbf{Monte Carlo (MC)}                        & \textbf{TD(0)}                                   \\ \hline
			\textbf{Requiere episodios completos} & Sí, requiere un episodio completo               & No, se actualiza paso a paso                     \\ \hline
			\textbf{Varianza}               & Alta, debido a las fluctuaciones del retorno     & Baja, ya que se basa en predicciones parciales    \\ \hline
			\textbf{Eficiencia}             & Puede ser menos eficiente en episodios largos    & Más eficiente, actualiza continuamente           \\ \hline
			\textbf{Estabilidad}            & Más estable en entornos con retorno claro       & Menos estable debido a la propagación de errores \\ \hline
			\textbf{Aplicabilidad}          & No es adecuado para entornos no episódicos       & Adecuado para entornos no episódicos             \\ \hline
		\end{tabularx}
		\end{minipage} 
	}}
	
    \item \key{Los métodos \textbf{n-step TD} generalizan} los métodos \textbf{MC} y \textbf{TD} para que uno pueda pasar suavemente de uno a otro según sea necesario para cumplir con las demandas de una tarea en particular.
    \item Los métodos \textbf{MC} están en un extremo del espectro y los métodos \textbf{TD de un solo paso} en el otro.
    \item \key{Los mejores métodos} suelen ser aquellos que se encuentran entre estos dos extremos.
    
    \begin{itemize}
    \item Conviene explorar en intervalo de tiempo más amplios.
    \end{itemize}
    
    \item \key{Estudiaremos} primero la predicción y después el control
\end{itemize}
\end{frame}


%%%%%%%%%%%%%%%%%%%%%%%%%%%%%%%%%%%%%%%%%%%%%%%%%%%%%%%%%%%%%%%%%%%%
%%%%%%%%%%%%%%%%%%%%%%%%%%%%%%%%%%%%%%%%%%%%%%%%%%%%%%%%%%%%%%%%%%%%

\section{Predicción TD de n pasos}


%%%%%%%%%%%%%%%%%%%%%%%%%%%%%%%%%%
%%%%%%%%%%%%%%%%%%%%%%%%%%%%%%%%%%
\begin{frame}{Predicción n-step TD}
\begin{itemize}
    \item Realizar una actualización basada en \key{una secuencia de \textbf{n} recompensas}:
    \begin{itemize}
        \item Más de una, pero menos que todas ellas.
    \end{itemize}
    \item \key{Gráficamente}, se pueden representan las siguientes variantes de TD:

	\centerline{\includegraphics[width=.5\textwidth]{fig/backup_diagrams_of_n_step}}

	\end{itemize}

\end{frame}

%%%%%%%%%%%%%%%%%%%%%%%%%%%%%%%%%%
%%%%%%%%%%%%%%%%%%%%%%%%%%%%%%%%%%

\begin{frame}{Predicción n-step TD (Actualización)}
\begin{itemize}
	\item Recuerda: $\color{violet} {Nuevo Estimado} = {Viejo Estimado} + {Tama\text{\itñ}o del Paso} \cdot \left( {Objetivo} - {Viejo Estimado} \right).$
	
    \item El \key{objetivo} (target) para una actualización arbitraria de n-step es el \key{retorno n-step}:
    \begin{align*}
    	G_{t:t+1} & \doteq  R_{t+1} + \gamma V_t(S_{t+1}) \\
    	G_{t:t+2} & \doteq  R_{t+1} + \gamma R_{t+2} + \gamma^2 V_{t+1}(S_{t+2}), \\
    	& \cdots  \\
		G_{t:t+n} & \doteq R_{t+1} + \gamma R_{t+2} + \dots + \gamma^{n-1} R_{t+n} + \gamma^{n} V_{t+n-1}(S_{t+n}) \quad \forall n\geq 1, \forall t \text{ t.q }  0 \leq t < T-n  \\
		G_{t:t+n} & \doteq G_t = R_{t+1} + \gamma R_{t+2} + \gamma^2 R_{t+3} + \dots + \gamma^{T-t-1} R_T \text{ si } t+n\geq T
	\end{align*} 

	
	\vskip -1.0cm
	
	~
    
    \item \key{Los retornos n-step} pueden considerarse \key{aproximaciones} del retorno completo, \key{truncadas} después de \(n\) pasos y luego corregidas por los términos faltantes.

    \item La fórmula para la \key{actualización n-step} es:
    $
    V_{t+n}(S_t) = V_{t+n-1}(S_t) + \alpha \left( {\color{red} G_{t:t+n}} - V_{t+n-1}(S_t) \right)
    $
    \item Mientras que los valores de todos los demás estados permanecen inalterados: 
    
     \hfil $
    V_{t+n}(s) = V_{t+n-1}(s) \quad \text{para todos } s \neq S_t
    $
    \item Una \key{expresión equivalente} es la siguiente: 
    
 
    \centerline{$
     V(S_t) \leftarrow V(S_t) + \alpha \left[ {\color{reddark} \sum_{k=1}^{n} \gamma^{k-1} R_{t+k} + \gamma^n V(S_{t+n})} - V(S_t) \right]
    $}
    
 
\end{itemize}
	
\end{frame}





%%%%%%%%%%%%%%%%%%%%%%%%%%%%%%%%%%
%%%%%%%%%%%%%%%%%%%%%%%%%%%%%%%%%%
\begin{frame}{DT de $n$-pasos para estimar $v_\pi$}{\cite{Sutton2018}}
	
	\centerline{\footnotesize$
	V(S_t) \leftarrow V(S_t) + \alpha \left[ \sum_{k=1}^{n} \gamma^{k-1} R_{t+k} + \gamma^n V(S_{t+n}) - V(S_t) \right]
	$}
	
	\begin{algorithm}[H]
		\centering
		\caption{DT de $n$-pasos para estimar $V\approx v_\pi$}
		\includegraphics[width=.67\textwidth]{fig/alg_n_step_TD}
	\end{algorithm}
	

\end{frame}


%%%%%%%%%%%%%%%%%%%%%%%%%%%%%%%%%%
%%%%%%%%%%%%%%%%%%%%%%%%%%%%%%%%%%

\begin{frame}{Un ejemplo para 3 pasos}{$G_{t:t+3} \doteq  R_{t+1} + \gamma R_{t+2} + \gamma^2 R_{t+3} + \gamma^3 V_{t+2}(S_{t+3})$}
\centerline{
\scalebox{.9}{
\begin{minipage}{.9\textwidth}
	
	
	\begin{tikzpicture}[->, >=Stealth, node distance=2cm, state/.style={circle, draw, minimum size=0.8cm}]
		\node[state] (S0) {S$_0$};
	\end{tikzpicture}
	
	

	\begin{tikzpicture}[->, >=Stealth, node distance=2cm, state/.style={circle, draw, minimum size=0.8cm}]
		
		% Nodos
		\node[state] (S0) {S$_0$};
		\node[state, right of=S0] (S1) {S$_1$};

		
		% Transiciones
		\draw (S0) edge node[above] {R$_1$} (S1);
		\node[right=0mm of S1] {\fbox{$t=0$}};
				
	\end{tikzpicture}
	
	\begin{tikzpicture}[->, >=Stealth, node distance=2cm, state/.style={circle, draw, minimum size=0.8cm}]
		
		% Nodos
		\node[state] (S0) {S$_0$};
		\node[state, right of=S0] (S1) {S$_1$};
		\node[state, right of=S1] (S2) {S$_2$};

		
		% Transiciones
		\draw (S0) edge node[above] {R$_1$} (S1);
		\draw (S1) edge node[above] {R$_2$} (S2);

		
		
	\end{tikzpicture}
	
	\begin{tikzpicture}[->, >=Stealth, node distance=2cm, state/.style={circle, draw, minimum size=0.8cm}]
		
		% Nodos
		\node[state] (S0) {$\color{red} S_0$};
		\node[state, right of=S0] (S1) {S$_1$};
		\node[state, right of=S1] (S2) {S$_2$};
		\node[state, right of=S2] (S3) {S$_3$};

		
		% Transiciones
		\draw (S0) edge node[above] {R$_1$} (S1);
		\draw (S1) edge node[above] {$\gamma R_2$} (S2);
		\draw (S2) edge node[above] {$\gamma^2 R_3$} (S3);
		
		\node[right=0mm of S3] (form3) {$\gamma^3 V(S_3)$};
		\node[right=5mm of form3] {\fbox{$t=2$} \fbox{$\tau=0$}};  
		
	\end{tikzpicture}
	
	\begin{tikzpicture}[->, >=Stealth, node distance=2cm, state/.style={circle, draw, minimum size=0.8cm}]
		
		% Nodos
		\node[state] (S0) {S$_0$};
		\node[state, right of=S0] (S1) {$\color{red} S_1$};
		\node[state, right of=S1] (S2) {S$_2$};
		\node[state, right of=S2] (S3) {S$_3$};
		\node[state, right of=S3] (S4) {S$_4$};
	
		
		% Transiciones
		\draw (S0) edge node[above] {} (S1);
		\draw (S1) edge node[above] {$R_2$} (S2);
		\draw (S2) edge node[above] {$\gamma R_3$} (S3);
		\draw (S3) edge node[above] {$\gamma^2 R_4$} (S4);
	
		\node[right=0mm of S4] (form4) {$\gamma^3 V(S_4)$};
		\node[right=5mm of form4] {\fbox{$t=3$} \fbox{$\tau=1$}}; 
		
	\end{tikzpicture}
	
	
	\begin{tikzpicture}[->, >=Stealth, node distance=2cm, state/.style={circle, draw, minimum size=0.8cm}]
		
		% Nodos
		\node[state] (S0) {S$_0$};
		\node[state, right of=S0] (S1) {S$_1$};
		\node[state, right of=S1] (S2) {$\color{red} S_2$};
		\node[state, right of=S2] (S3) {S$_3$};
		\node[state, right of=S3] (S4) {S$_4$};
		\node[state, right of=S4] (S5) {S$_5$};
		
		% Transiciones
		\draw (S0) edge node[above] {} (S1);
		\draw (S1) edge node[above] {} (S2);
		\draw (S2) edge node[above] {$R_3$} (S3);
		\draw (S3) edge node[above] {$\gamma R_4$} (S4);
		\draw (S4) edge node[above] {$\gamma^2 R_5$} (S5);
		\node[right=0mm of S5] (T5) {\fbox{$T=5$}};
		\node[right=0mm of T5] {\fbox{$t=4$} \fbox{$\tau=2$}};
		
	\end{tikzpicture}
	
		\begin{tikzpicture}[->, >=Stealth, node distance=2cm, state/.style={circle, draw, minimum size=0.8cm}]
		
		% Nodos
		\node[state] (S0) {S$_0$};
		\node[state, right of=S0] (S1) {S$_1$};
		\node[state, right of=S1] (S2) {$S_2$};
		\node[state, right of=S2] (S3) {$\color{red}  S_3$};
		\node[state, right of=S3] (S4) {S$_4$};
		\node[state, right of=S4] (S5) {S$_5$};
		
		% Transiciones
		\draw (S0) edge node[above] {} (S1);
		\draw (S1) edge node[above] {} (S2);
		\draw (S2) edge node[above] {} (S3);
		\draw (S3) edge node[above] {$R_4$} (S4);
		\draw (S4) edge node[above] {$\gamma R_5$} (S5);
		
		\node[right=0mm of S5] {\fbox{$t=5$} \fbox{$\tau=3$}};
	\end{tikzpicture}
	
	
	
\begin{tikzpicture}[->, >=Stealth, node distance=2cm, state/.style={circle, draw, minimum size=0.8cm}]
	
	% Nodos
	\node[state] (S0) {S$_0$};
	\node[state, right of=S0] (S1) {S$_1$};
	\node[state, right of=S1] (S2) {S$_2$};
	\node[state, right of=S2] (S3) {S$_3$};
	\node[state, right of=S3] (S4) {$\color{red}  S_4$};
	\node[state, right of=S4] (S5) {S$_5$};
	
	% Transiciones
	\draw (S0) edge node[above] {} (S1);
	\draw (S1) edge node[above] {} (S2);
	\draw (S2) edge node[above] {} (S3);
	\draw (S3) edge node[above] {} (S4);
	\draw (S4) edge node[above] {$R_5$} (S5);
	
	\node[right=0mm of S5] {\fbox{$t=6$} \fbox{$\tau=4$}};
	
\end{tikzpicture}
\end{minipage}}}


\TPMargin{0mm}%
\begin{textblock}{9}(7, 3) \small
\includegraphics[width=\textwidth,trim=35 5 5 350, clip]{fig/alg_n_step_TD}
\end{textblock}


\end{frame}





%%%%%%%%%%%%%%%%%%%%%%%%%%%%%%%%%%
%%%%%%%%%%%%%%%%%%%%%%%%%%%%%%%%%%
	
\begin{frame}{Propiedades de los Retornos de \( n \) Pasos en TD}
	
	\begin{itemize}
		\item \textbf{Retornos de \( n \) pasos} utilizan la función de valor \( V_{t+n-1} \) para corregir las recompensas faltantes más allá de \( R_{t+n} \).
		
		\item El retorno de \( n \) pasos toma en cuenta más información que solo el valor \( V_{t+n-1}(s) \), ya que incluye las recompensas obtenidas durante los \( n \) pasos y una estimación de los valores futuros (a través de \( V_{t+n-1}(s_{t+n}) \)).

		\item La propiedad de reducción de error garantiza que el peor error del retorno esperado de \( n \) pasos es:
		\[\ds
		\max_s \Big| E_{\pi}[G_{t:t+n} | S_t = s] - v_{\pi}(s) \Big| \leq \gamma^n \max_s \Big| V_{t+n-1}(s) - v_{\pi}(s) \Big|
		\]
		\item \key{Convergencia}: Los métodos TD de \( n \) pasos convergen a las predicciones correctas bajo condiciones técnicas adecuadas.
		\item Los métodos TD de \( n \) pasos son una \textbf{familia sólida de métodos}, con los métodos TD de un solo paso y Monte Carlo como miembros extremos.
		
		\centerline{\textbf{Pero depende de $\alpha$}}
	\end{itemize}

\end{frame}
	



%%%%%%%%%%%%%%%%%%%%%%%%%%%%%%%%%%
%%%%%%%%%%%%%%%%%%%%%%%%%%%%%%%%%%

%%%%%%%%%%%%%%%%%
%%%%%%%%%%%%%%%%%
\begin{frame}{Ejemplo: ¿Qué valor de $n$ es mejor?}{$n$ pasos en el Random Walk}
	\centerline{
		\includegraphics[width=0.8\textwidth]{fig/randomWalk}
	}
	
	\begin{columns}
		\begin{column}{.5\textwidth}
			\includegraphics[width=\textwidth]{fig/Performance_of_n_step_TD}
		\end{column}
		%
		\begin{column}{.5\textwidth}
			\begin{itemize}
				\item 19 estados, saliendo por la izquierda con -1.
				
				\item Rendimiento de los métodos TD de $n$ pasos en función de \( \gamma \), para varios valores de \( n \), en una tarea de random walk de 19 estados.
				
				\item La medida de rendimiento  es la raíz cuadrada del error cuadrada promedio entre las predicciones al final del episodio para los 19 estados y sus valores verdaderos, entonces se promedian sobre los primeros 10 episodios y se realizan 100 repeticiones de todo el experimento (se usaron los mismos conjuntos de caminatas para todas las configuraciones de parámetros). 
				
				
			\end{itemize}
		\end{column}
	\end{columns}
\end{frame}





%%%%%%%%%%%%%%%%%%%%%%%%%%%%%%%%%%%%%%%%%%%%%%%%%%%%%%%%%%%%%%%%%%%%
%%%%%%%%%%%%%%%%%%%%%%%%%%%%%%%%%%%%%%%%%%%%%%%%%%%%%%%%%%%%%%%%%%%%

\section{Control TD de n pasos}



%%%%%%%%%%%%%%%%%%%%%%%%%%%%%%%%%%%%%%%%%%%%%%%%%%%%%%%%%%%%%%%%%%%%
%%%%%%%%%%%%%%%%%%%%%%%%%%%%%%%%%%%%%%%%%%%%%%%%%%%%%%%%%%%%%%%%%%%%
\subsection{Aprendizaje On-Policy de $n$ Pasos}



%%%%%%%%%%%%%%%%%%%%%%%%%%%%%%%%%%
%%%%%%%%%%%%%%%%%%%%%%%%%%%%%%%%%%



%%%%%%%%%%%%%%%%%%%%%%%%%%%%%%%%%%
%%%%%%%%%%%%%%%%%%%%%%%%%%%%%%%%%%
\begin{frame}{Variantes de SARSA}
	\begin{itemize}
		\item Realizar una actualización basada en una \key{secuencia de \textbf{n} recompensas}:
		\begin{itemize}
			\item Más de una, pero menos que todas ellas.
		\end{itemize}
		\item \key{Gráficamente}, se representan las siguientes variantes de Sarsa:
		
		\centerline{\includegraphics[width=.7\textwidth]{fig/backup_diagrams_of_n_step_sarsa}}
		
	\end{itemize}
	
	
	\note<1>{
		\begin{itemize}
			\item Se muestran los Diagramas de respaldo para los métodos de \(n\)-pasos para valores estado–acción. Estos métodos abarcan desde la actualización de un paso de Sarsa(0) hasta la actualización de todo el episodio del método Monte Carlo. 
			
			\item En el medio se encuentran las actualizaciones de \(n\)-pasos, basadas en \(n\) pasos de recompensas reales y el valor estimado del \(n\)-ésimo siguiente par estado–acción, todos descontados adecuadamente. 
			
			\item En el extremo derecho se encuentra el diagrama de respaldo para el método de \(n\)-pasos de Sarsa Esperado.
			
			\item Consiste en una cadena lineal de acciones y estados de muestra, igual que en n-pasos Sarsa, excepto que su último elemento es una rama sobre todas las posibilidades de acción ponderadas por su probabilidad bajo la política \( \pi \).
			
			\item Vamos mostrar cómo se actualizan las funciones de valor sobre las acciones tanta para el Sarsa de n-pasos como el Sarsa esperado de n-pasos.
		\end{itemize}
		
	\fin
	}
\end{frame}
			
	
			

%%%%%%%%%%%%%%%%%
%%%%%%%%%%%%%%%%%
\begin{frame}{Sarsa's de $n$-pasos}
	

\begin{itemize}
	\item  La \key{versión original} es Sarsa de un solo paso, o Sarsa(0). 
	
	\centerline{ $ Q(S_t, A_t) \gets Q(\underline{S_t, A_t}) + \alpha \left( \underline{{\color{red} R_{t+1}}} {\color{red} + \gamma Q (}\underline{{\color{red}S_{t+1}, A_{t+1})}} - Q(S_t, A_t) \right) $}
	
	
	\item  La \textit{versión de n pasos} de Sarsa la llamamos \key{n-step Sarsa}.
	
	\begin{itemize}
		\item  El \key{objetivo} (target) para una actualización arbitraria de n-step es el \key{retorno n-step}:
		\begin{align*}
			G_{t:t+1} & \doteq  R_{t+1} + \gamma Q_t(S_{t+1}, A_{t+1}) \\
			G_{t:t+2} & \doteq  R_{t+1} + \gamma R_{t+2} + \gamma^2 Q_{t+1}(S_{t+2}, A_{t+2}) \\
			& \cdots  \\
			{\color{red} G_{t:t+n} } & \doteq R_{t+1} + \gamma R_{t+2} + \dots + \gamma^{n-1} R_{t+n} + \gamma^n {\color{violet} Q_{t+n-1} (S_{t+n}, A_{t+n})},	\text{ para }  n \geq 1 \text{ y } 0 \leq t < T  \\
			G_{t:t+n} & \doteq G_t = R_{t+1} + \gamma R_{t+2} + \gamma^2 R_{t+3} + \dots + \gamma^{T-t-1} R_T \text{ si } t+n\geq T
		\end{align*} 
		
		\vspace{-0.5cm}
	
	
	
		\item La \key{fórmula de actualización} es: 
			{\normalsize $Q_{t+n}(S_t, A_t) = Q_{t+n-1}(S_t, A_t) + \alpha \big( {\color{red} G_{t:t+n}} - Q_{t+n-1}(S_t, A_t) \big)$}
	
	\end{itemize}
	
	\item \key{n-step Expected SARSA} es la versión de n pasos de Expected Sarsa
		
		\begin{itemize}
			
			\item  Se redefine el \key{objetivo} (target), llamado \key{retorno n-step}, y se usa fórmula de actualización anterior:
			
			{\normalsize $ {\color{red}G_{t:t+n}}  \doteq R_{t+1} + \cdots + \gamma^{n-1} R_{t+n} + \gamma^n {\color{violet} \overline{V}_{t+n-1} (S_{t+n})}, \; t+n < T $}
			\hfil y \hfil
			$ G_{t:t+n}  \doteq G_t  \text{ para } t+n \geq T$

				
			\item[] donde $\ds \overline{V}_{t}(s) \doteq \sum_a \pi(a|s) Q_t(s, a), \quad \forall s \in S.$
				
		\end{itemize}
		
\end{itemize}

	
\end{frame}





%%%%%%%%%%%%%%%%%%%%%%%%%%%%%%%%%%
%%%%%%%%%%%%%%%%%%%%%%%%%%%%%%%%%%
\begin{frame}{SARSA de $n$-pasos para estimar $Q\approx q_\pi$}{$Q_{t+n}(S_t, A_t) = Q_{t+n-1}(S_t, A_t) + \alpha \left( {\color{red}G_{t:t+n}} - Q_{t+n-1}(S_t, A_t) \right)$}

		\centering
		\includegraphics[width=.67\textwidth]{fig/alg_n_step_sarsa}
		%
		\raisebox{1cm}{
		
	\begin{minipage}{.23\textwidth}\small 
	Convergencia de \( Q \):
	
	\footnotesize
	\mbox{$
	\max_s \max_a | Q_{\text{new}}(s, a) - Q_{\text{old}}(s, a) | < \delta
	$}
	
	$
	\text{max}_s |\pi_{\text{new}}(s) - \pi_{\text{old}}(s)| < \delta
	$
	\end{minipage}}

\end{frame}



%%%%%%%%%%%%%%%%%
%%%%%%%%%%%%%%%%%
\begin{frame}{Ejemplo de SARSA de $n$-pasos}
	
	\includegraphics[width=\textwidth]{fig/gridWorld_n_steps_sarsa}
	
	\begin{itemize}
		\item \key{Izquierda:} Un episodio.
		\item \key{Centro:} SARSA a 1 paso. Solo la última acción de la secuencia de acciones actualiza la recompensa.
		\item \key{Derecha:} SARSA a 10 pasos. Las últimas 10 acciones de la secuencia de acciones actualizan sus valores. Aprende mucho más rápido en un solo episodio.
	\end{itemize}

\end{frame}


%%%%%%%%%%%%%%%%%%%%%%%%%%%%%%%%%%%%%%%%%%%%%%%%%%%%%%%%%%%%%%%%%%%%
%%%%%%%%%%%%%%%%%%%%%%%%%%%%%%%%%%%%%%%%%%%%%%%%%%%%%%%%%%%%%%%%%%%%

\subsection{Aprendizaje Off-Policy de $n$ Pasos}



%%%%%%%%%%%%%%%%%%%%%%%%%%%%%%%%%%%%%%%%%%%%%%%%%%%%%%%%%%%%%%%%%%%%
%%%%%%%%%%%%%%%%%%%%%%%%%%%%%%%%%%%%%%%%%%%%%%%%%%%%%%%%%%%%%%%%%%%%

\subsubsection{Basado en Muestreo por Importancia}

%%%%%%%%%%%%%%%%%%%%%%%%%%%%%%%%%%
%%%%%%%%%%%%%%%%%%%%%%%%%%%%%%%%%%
\begin{frame}[label={frame:RecordatorioMC}]{Recordatorio MC}
	\begin{itemize}
		\item En \key{Monte Carlo} estudiamos el muestreo por importancia ordinario y ponderado. 
		
		\centerline{$\ds 
		V(s) \doteq \frac{\sum_{t \in \mathcal{T}(s)} \rho_{t:T(t)-1} G_t}{|\mathcal{T}(s)|}
		$
		\qquad
		$\ds 
		V(s) \doteq \frac{\sum_{t \in \mathcal{T}(s)} \rho_{t:T(t)-1} G_t}{\sum_{t \in \mathcal{T}(s)} \rho_{t:T(t)-1}},
		$
		\qquad con $\ds \rho_{t:T(t)-1} \doteq \prod_{k=t}^{T(t)-1} \frac{\pi(A_k | S_k)}{b(A_k | S_k)}$}
		
		\item En su momento vimos que las funciones $Q(s,a)$ tienen  un \key{cálculo incremental}:
		\[  Q(S_t, A_t) \gets Q(S_t, A_t) + \frac{\rho_{t:T(t)-1}}{C(S_t, A_t)}   [G_t -  Q(S_t, A_t)]
		= Q(S_t, A_t) + \alpha \rho_{t:T(t)-1}  [G_t -  Q(S_t, A_t)] \]		

		\item \key{En métodos de $n$ pasos}, estamos interesados en 
		
		\begin{itemize}
			\item La ganancia para n-pasos: 
				\begin{itemize}
					\item $\ds G_{t:t+n} = \underbrace{R_{t+1} + \gamma R_{t+2} + \dots + \gamma^{n-1} R_{t+n}}_{(n-1)-recompensas} + \underbrace{\gamma^n FuncionValor}_{Recompensa-Truncada}$, para \( n \geq 1 \) y \( 0 \leq t +n < T \)
					
					\item $\ds G_{t:t+n} = G_{t:T}$ si $t+n\geq T$
				\end{itemize}
				
			
			\item la probabilidad relativa de solo  $n$ acciones. $\ds \rho_{t:\lim} = \prod_{k=t}^{\lim} \frac{\pi(A_k | S_k)}{b(A_k | S_k)}$
		\end{itemize}
	\end{itemize}

\end{frame}







%%%%%%%%%%%%%%%%%%%%%%%%%%%%%%%%%%
%%%%%%%%%%%%%%%%%%%%%%%%%%%%%%%%%%
\begin{frame}{Aprendizaje Off-Policy de $n$ Pasos}
	\begin{itemize}
		\item Definiremos la \key{razón de muestreo de importancia} de tomar las $n$-acciones (desde $t$ hasta $h$):
		\centerline{$\ds
			\rho_{t:h} \doteq \prod_{k=t}^{\min(h,T-1)} \frac{\pi(A_k | S_k)}{b(A_k | S_k)}
			$}
			
			\begin{itemize}
				\item Donde $T$ hace referencia al instante de llegar al estado final.
			\end{itemize}
			
			
		\item \key{Para valores de estado}, definiremos la función de actualización como
		
		\centerline{\normalsize
			$V_{t+n}(S_t) \doteq V_{t+n-1}(S_t) + \alpha$ \Large $ \color{reddark}\rho_{t:t+n-1} $ \normalsize $\big[G_{t:t+n} - V_{t+n-1}(S_t)\big] \quad 0 \leq t < T$ }
			
		\begin{itemize}
			\item $ {\color{red}G_{t:t+n}}  \doteq R_{t+1} + \gamma R_{t+2} + \dots + \gamma^{n-1} R_{t+n} + \gamma^{n} {\color{violet}  V_{t+n-1}(S_{t+n})}$
			\item La {razón de muestreo de importancia} de tomar las $n$ acciones va desde $A_t$ hasta $A_{t+n-1}$
		\end{itemize}



 
 

		\item \key{Para valores de acción}, definiremos la función de actualización como
		 
		\centerline{\normalsize
			$Q_{t+n}(S_t, A_t) \doteq Q_{t+n-1}(S_t, A_t) + \alpha$ \Large $\color{reddark} \rho_{t+1:t+n} $ \normalsize $\big[ G_{t:t+n} - Q_{t+n-1}(S_t, A_t) \big], \quad 0 \leq t < T,$ }
		
		\begin{itemize}			
			\item $ {\color{red}G_{t:t+n}}  \doteq R_{t+1} + \gamma R_{t+2} + \dots + \gamma^{n-1} R_{t+n} + \gamma^n {\color{violet} Q_{t+n-1} (S_{t+n}, A_{t+n})}$
			
			\item La {razón de muestreo de importancia} de tomar las $n$ acciones va desde $A_t$ hasta $A_{t+n-1}$
			
		\end{itemize}
				
	
		\item \key{Si las dos políticas son las mismas}, el ratio de muestreo por importancia es 1.
		\begin{itemize}
			\item $V_{t+n}(S_t)$  generaliza la actualización TD de n-pasos.
			\item $Q_{t+n}(S_t, A_t)$  generaliza la actualización Sarsa de n-pasos.
		\end{itemize} 
		
		\item Para la versión off-policy de \key{n-step Expected Sarsa}:
			\begin{itemize}
				\item $ {\color{red}G_{t:t+n}}  \doteq R_{t+1} + \cdots + \gamma^{n-1} R_{t+n} + \gamma^n {\color{violet} \overline{V}_{t+n-1} (S_{t+n})}$ \quad donde  $\overline{V}_{t}(s) \doteq \sum_a \pi(a|s) Q_t(s, a)$. 
				\item Usará la función $Q$ con {\Large $\color{reddark} \rho_{t+1:t+n-1} $} en lugar de \( \rho_{t+1:t+n} \)
				
			\end{itemize}
		
%		La versión off-policy de n-pasos Expected Sarsa usaría la misma actualización que la anterior para n-pasos Sarsa, excepto que la razón de muestreo de importancia tendría un factor menos en ella. Es decir, la ecuación anterior usaría \( \rho_{t+1:t+n-1} \) en lugar de \( \rho_{t+1:t+n} \), y, por supuesto, usaría la versión Expected Sarsa del retorno n-pasos (7.7). Esto se debe a que en Expected Sarsa todas las acciones posibles se tienen en cuenta en el último estado; la que realmente se toma no tiene efecto y no es necesario corregirla.
		
		
		
		
		
	\end{itemize}
	
\end{frame}




%%%%%%%%%%%%%%%%%%%%%%%%%%%%%%%%%%
%%%%%%%%%%%%%%%%%%%%%%%%%%%%%%%%%%
\begin{frame}{SARSA de $n$-pasos off-policy para estimar $Q\approx q_\pi$}
	{\small $Q_{t+n}(S_t, A_t) \doteq Q_{t+n-1}(S_t, A_t) + \alpha  \rho_{t+1:t+n} \big[ G_{t:t+n} - Q_{t+n-1}(S_t, A_t) \big]$}
	
	\centering
	\includegraphics[width=.66\textwidth]{fig/alg_off_policy_n_step_sarsa}	

\end{frame}



%%%%%%%%%%%%%%%%%%%%%%%%%%%%%%%%%%%%%%%%%%%%%%%%%%%%%%%%%%%%%%%%%%%%
%%%%%%%%%%%%%%%%%%%%%%%%%%%%%%%%%%%%%%%%%%%%%%%%%%%%%%%%%%%%%%%%%%%%

\subsubsection{Sin Muestreo por Importancia}




%%%%%%%%%%%%%%%%%
%%%%%%%%%%%%%%%%%
\begin{frame}{Idea Intuitiva de \textit{n-step Tree Backup Algorithm}}
	\begin{columns}
		\begin{column}{.87\textwidth}
			\begin{itemize}
				\item ¿Es posible el aprendizaje off-policy sin muestreo de importancia?
				\item Árbol de respaldo para 3 acciones:
				\begin{itemize}
					\item Hay 3 muestras de estados y recompensas, y 2 muestras de acciones
					\item  Las v.a. representan los \key{eventos que ocurren} \textbf{después} del par inicial estado-acción \( (S_t, A_t) \).
					\item Colgadas a los lados de cada estado están las \key{acciones que no fueron seleccionadas}. 
					\item Para el último estado, todas las acciones se consideran como \textit{acciones que no han sido (todavía) seleccionadas}.
					\begin{itemize}
						\item Debido a que no tenemos datos de muestra para las acciones no seleccionadas, realizamos \textit{bootstrap} y utilizamos las estimaciones de sus valores para formar el objetivo para la actualización.
					\end{itemize}
				\end{itemize}
				
				\item  Cada nodo hoja contribuye al objetivo con un peso proporcional a su probabilidad de ocurrir bajo la política objetivo \( \pi \).
				\begin{itemize}
					
					\item Cada acción de primer nivel \( a \) contribuye con un peso de \( \pi(a|S_{t+1}) \).
					\item La acción realmente tomada, \( A_t \), no contribuye en absoluto. Su probabilidad, \( \pi(A_{t+1}|S_{t+1}) \), se utiliza para ponderar todos los valores de acción del segundo nivel.
					\item Cada acción de segundo nivel no seleccionada, \( a' \), contribuye con peso \( \pi(A_{t+1}|S_{t+1}) \pi(a'|S_{t+2}) \).
					\item Cada acción de tercer nivel contribuye con peso \( \pi(A_{t+1}|S_{t+1}) \pi(A_{t+2}|S_{t+2}) \pi(a''|S_{t+3}) \), y así sucesivamente.
				\end{itemize}
			\end{itemize}
		\end{column}
		%
		\begin{column}{.12\textwidth}
			\includegraphics[width=\textwidth]{fig/tree_backup_3_step}
		\end{column}
	\end{columns}
\end{frame}




%%%%%%%%%%%%%%%%%
%%%%%%%%%%%%%%%%%
\begin{frame}{Justificación Matemática}

	\begin{itemize}
		\item El retorno de un solo paso (objetivo) lo vamos a considera igual  que  Expected Sarsa:
		
		\centerline{$\ds G_{t:t+1} = R_{t+1} + \gamma \sum_a \pi(a|S_{t+1}) Q_{t}(S_{t+1}, a) \text{ si }  t < T-1$ ;\qquad $G_{T-1:T} = R_{T}$ \text{ si }  t = T-1}
		
		\item El retorno de dos pasos es para \( t < T-2 \)
		
		\vspace{-0.5cm}
		
		{\small
		\begin{align*}
			G_{t:t+2} &\doteq R_{t+1} + \gamma {\color{reddark} \sum_{a \neq A_{t+1}} \pi(a \mid S_{t+1}) Q_{t+1}(S_{t+1}, a)}
			    + \gamma {\color{ForestGreen} \pi(A_{t+1} \mid S_{t+1}) }  \left( R_{t+2} + \gamma \sum_{a} \pi(a \mid S_{t+2}) Q_{t+1}(S_{t+2}, a) \right) \\
			&= R_{t+1} + \gamma {\color{reddark} \sum_{a \neq A_{t+1}} \pi(a \mid S_{t+1}) Q_{t+1}(S_{t+1}, a)}
			 + \gamma {\color{ForestGreen} \pi(A_{t+1} \mid S_{t+1}) } {\color{black} G_{t+1:t+\mbf{2}}},
		\end{align*}
		}
		
		\item El retorno para  n-pasos, $t < T-1,  n \geq 2$:
		
		\centerline{$\ds G_{t:t+n} \doteq R_{t+1} + \gamma {\color{reddark} \sum_{a \neq A_{t+1}} \pi(a|S_{t+1}) Q_{t+n-1}(S_{t+1}, a)} + \gamma {\color{ForestGreen} \pi(A_{t+1}|S_{t+1}) } {\color{black} G_{t+1:t+\mbf{n}}}$}
		
		\item Regla de actualización de valor-acción de n-pasos Sarsa:
		
		\centerline{\large $\ds	Q_{t+n}(S_t, A_t) \doteq Q_{t+n-1}(S_t, A_t) + \alpha \left[ G_{t:t+n} - Q_{t+n-1}(S_t, A_t) \right], $}\vspace{0.25cm}
		
		para \( 0 \leq t < T \), mientras que los valores de todos los demás pares estado-acción permanecen sin cambios: \( Q_{t+n}(s, a) = Q_{t+n-1}(s, a) \), para todos \( s, a \) tales que \( s \neq S_t \) o \( a \neq A_t \).
		
	\end{itemize}
	
\end{frame}




%%%%%%%%%%%%%%%%%%%%%%%%%%%%%%%%%%
%%%%%%%%%%%%%%%%%%%%%%%%%%%%%%%%%%
\begin{frame}{Árbol de Respaldo de $n$-pasos para estimar $Q\approx q_\pi$}
	{\small $Q_{t+n}(S_t, A_t) \doteq Q_{t+n-1}(S_t, A_t) + \alpha  \big[ G_{t:t+n} - Q_{t+n-1}(S_t, A_t) \big]$}
	
	
	
	
	\begin{columns}
		\begin{column}{.87\textwidth}
			\centering
			\includegraphics[width=.7\textwidth]{fig/alg_n_step_tree_backup}	
		\end{column}
		%
		\begin{column}{.10\textwidth}
			\includegraphics[width=\textwidth]{fig/tree_backup_3_step}
		\end{column}
	\end{columns}

\end{frame}






%%%%%%%%%%%%%%%%%%%%%%%%%%%%%%%%%%
%%%%%%%%%%%%%%%%%%%%%%%%%%%%%%%%%%
% BIBLIOGRAFÍA
%%%%%%%%%%%%%%%%%%%%%%%%%%%%%%%%%%
%%%%%%%%%%%%%%%%%%%%%%%%%%%%%%%%%%


%%%%%%%%%%%%%%%%%%%%%%%%%%%%%%%%%%%%%%%%%%%%%%%%%%%%%%%%%%%%%%%%%%%%
%%%%%%%%%%%%%%%%%%%%%%%%%%%%%%%%%%%%%%%%%%%%%%%%%%%%%%%%%%%%%%%%%%%%

\section{Referencias}



% Redefinimos la apariencia
%\setbeamertemplate{headline}{}
%\setbeamertemplate{frametitle}{\vskip 0.5cm	\small{\textbf{\insertframetitle}}}
%\setbeamertemplate{footline}{}






% Mostramos la bibliografía
% \begin{frame}[allowframebreaks] %% Si ocupa más de una página
	%\begin{frame}[allowframebreaks]{Referencias -} %% Si hay sólo una página
	
	\begin{frame}{Referencias} %% Si hay sólo una página
		\nocite{Sutton2018, silver2015, Hado2018Lecture3}
		
		% \vskip -5cm %% Para ajustar la distancia al comienzo	
		%\bibliography{references}	
		\printbibliography[heading=none]
	\end{frame}    
	
	
	
	%%%%%%%%%%%%%%%%%%%%%%%%%%%%%%%%%%
	%%%%%%%%%%%%%%%%%%%%%%%%%%%%%%%%%%
	


%%%%%%%%%%%%%%%%%%%%%%%%%%%%%%%%%%
%%%%%%%%%%%%%%%%%%%%%%%%%%%%%%%%%%


\end{document}



